\documentclass[11pt,a4paper]{article}
\usepackage[utf8]{inputenc}
\usepackage[margin=1in]{geometry}
\usepackage{hyperref}
\usepackage{enumitem}
\usepackage{xcolor}
\usepackage{booktabs} % Added for better table formatting

% Color setup for links
\hypersetup{
    colorlinks=true,
    linkcolor=blue,
    filecolor=magenta,      
    urlcolor=cyan,
}

% ==============================================================================
% CUSTOM AI-FILLABLE MACRO
% ==============================================================================
% Usage:
% \papersummary{Title}{Authors}{Year/Venue}{Core Problem}{Architecture/Method}{Hardware/Setup}{Results}{Results Analysis & Tables}{Limitations & Relevance}
% ==============================================================================

\newcommand{\papersummary}[9]{
    \vspace{1em}
    \noindent\fbox{%
    \begin{minipage}{\dimexpr\textwidth-2\fboxsep-2\fboxrule\relax}
        \vspace{0.5em}
        \begin{center}
            \Large \textbf{#1}
        \end{center}
        \vspace{0.2em}
        \noindent\textbf{Authors:} #2 \hfill \textbf{Year \& Venue:} #3 \par\vspace{0.5em}
        \noindent\textbf{Core Problem:} #4 \par\vspace{0.5em}
        \noindent\textbf{Architecture / Method:} #5 \par\vspace{0.5em}
        \noindent\textbf{Hardware, Sensors \& Datasets:} #6 \par\vspace{0.5em}
        \noindent\textbf{Key Results:} #7 \par\vspace{0.5em}
        \noindent\textbf{Results Analysis \& Tables:} \par\vspace{0.2em} 
        \noindent #8 \par\vspace{0.5em}
        \noindent\textbf{Relevance / Limitations:} #9
        \vspace{0.5em}
    \end{minipage}%
    }
    \vspace{1em}
}

\title{Literature Review Documentation}
\author{Research Log}
\date{\today}

\begin{document}

\maketitle

\section*{Overview}
This document aggregates literature summaries. An AI assistant can be instructed to read a paper and output exactly one \texttt{\textbackslash papersummary} macro block per paper.

\section*{Paper Summaries}

% --- PASTE AI OUTPUT BELOW THIS LINE ---

\papersummary
{Example Paper: Hierarchical Frameworks for Robotic Manipulation}
{Doe, J., Smith, A., et al.}
{2026, IEEE/RSJ IROS}
{Addressing the latency and generalization gap in fine-motor robotic tasks using standard monolithic models.}
{A Hierarchical Mixture of Experts (HMoE) architecture separating high-level task planning from low-level joint parsing.}
{Evaluated on a simulated 6-DoF arm with RGB-D camera inputs (similar to RealSense D435if). Uses custom action-chunking datasets.}
{Reduced inference latency by 15\% while increasing successful grasp rates on unseen geometries by 22\%.}
{
    The model demonstrated significant performance improvements in multi-stage manipulation tasks compared to the baseline.
    \vspace{0.5em}
    
    \begin{center}
    \begin{tabular}{lcc}
        \toprule
        \textbf{Method} & \textbf{Latency (ms)} & \textbf{Success Rate (\%)} \\
        \midrule
        Baseline Monolithic & 142 & 68.5 \\
        Proposed HMoE & \textbf{120} & \textbf{90.5} \\
        \bottomrule
    \end{tabular}
    \end{center}
}
{The action chunking size is fixed, which limits dynamic reactive adjustments. Highly relevant for structuring the modular components of our current pipeline.}

% Blank template for AI reference
% \papersummary
% {Title}
% {Authors}
% {Year, Venue}
% {Problem description.}
% {Methodology description.}
% {Hardware details.}
% {Results.}
% {Results Analysis and Tables.}
% {Relevance.}

\end{document}